\exercisesheader{}

% 1

\eoce{\qt{Identifikasi parameter, Bagian I\label{identify_parameter_1}} Untuk setiap situasi berikut, tentukan apakah parameter yang menjadi perhatian berupa rata-rata atau proporsi. Akan membantu jika Anda memperhatikan apakah tanggapan setiap individu bersifat numerik atau kategoris.
\begin{parts}
\item Dalam sebuah survei, seratus mahasiswa ditanya berapa jam per minggu yang mereka habiskan di internet.
\item Dalam sebuah survei, seratus mahasiswa ditanya: ``Berapa persen dari waktu yang Anda habiskan di internet digunakan untuk tugas kuliah?''
\item Dalam sebuah survei, seratus mahasiswa ditanya apakah mereka mengutip informasi dari Wikipedia dalam makalah mereka.
\item Dalam sebuah survei, seratus mahasiswa ditanya berapa persen dari seluruh pengeluaran mingguan mereka yang digunakan untuk minuman beralkohol.
\item Dalam sampel yang terdiri atas seratus lulusan perguruan tinggi baru, ditemukan bahwa 85 persen memperkirakan akan memperoleh pekerjaan dalam satu tahun setelah tanggal kelulusan.
\end{parts}
}{}

% 2

\eoce{\qt{Identifikasi parameter, Bagian II\label{identify_parameter_2}} Untuk setiap situasi berikut, tentukan apakah parameter yang menjadi perhatian berupa rata-rata atau proporsi.
\begin{parts}
\item Sebuah jajak pendapat menunjukkan bahwa 64\% warga Amerika secara pribadi sangat mengkhawatirkan belanja pemerintah federal dan defisit anggaran.
\item Sebuah survei melaporkan bahwa pendapatan dari siaran berita televisi lokal meningkat 17\% dalam kurun dua tahun, sedangkan pendapatan surat kabar menurun 6,4\% selama kurun tersebut.
\item Dalam sebuah survei, siswa sekolah menengah dan mahasiswa ditanya apakah mereka menggunakan layanan geolokasi pada ponsel pintar.
\item Dalam sebuah survei, pengguna ponsel pintar ditanya apakah mereka menggunakan layanan taksi berbasis web.
\item Dalam sebuah survei, pengguna ponsel pintar ditanya berapa kali mereka menggunakan layanan taksi berbasis web selama setahun terakhir.
\end{parts}
}{}

% 3

\eoce{\qt{Pengendalian mutu\label{comp_chips_quality_ctrl_prop}}
Sebagai bagian dari proses pengendalian mutu cip komputer,
seorang insinyur di sebuah pabrik mengambil sampel acak 212 cip
selama satu minggu produksi guna mengetahui proporsi cip
yang mengalami cacat berat pada saat itu.
Ia menemukan bahwa 27 cip mengalami cacat.
\begin{parts}
\item
    Populasi apa yang dipertimbangkan dalam kumpulan data ini?
\item
    Parameter apa yang sedang diestimasi?
\item\label{comp_chips_quality_ctrl_prop_pt_est}%
    Berapakah estimasi titik untuk parameter tersebut?
\item\label{comp_chips_quality_ctrl_prop_se_name}%
    Apa nama statistik yang kita gunakan untuk mengukur
    ketidakpastian estimasi titik?
\item\label{comp_chips_quality_ctrl_prop_se_calc_w_pt_est}%
    Hitung nilai dari
    bagian~(\ref{comp_chips_quality_ctrl_prop_se_name})
    untuk konteks ini.
\item
    Tingkat cacat historis adalah 10\%.
    Apakah insinyur itu seharusnya terkejut oleh tingkat cacat
    yang diamati selama minggu berjalan?
\item
    Misalkan nilai populasi yang sebenarnya diketahui sebesar 10\%.
    Jika kita menggunakan proporsi ini untuk menghitung ulang nilai pada
    bagian~(\ref{comp_chips_quality_ctrl_prop_se_calc_w_pt_est})
    dengan menggunakan $p = 0.1$ sebagai pengganti $\hat{p}$,
    apakah nilai yang dihasilkan banyak berubah?
\end{parts}
}{}

% 4

\eoce{\qt{Pengeluaran tak terduga\label{us_emergency_expense_prop}}
Dalam sampel acak yang terdiri atas 765 orang dewasa di Amerika Serikat, 322 orang mengatakan
bahwa mereka tidak mampu menutup pengeluaran tak terduga sebesar \$400 tanpa meminjam
uang atau berutang.
% Ref: https://www.federalreserve.gov/publications/files/2017-report-economic-well-being-us-households-201805.pdf
\begin{parts}
\item
    Populasi apa yang dipertimbangkan dalam kumpulan data ini?
\item
    Parameter apa yang sedang diestimasi?
\item\label{us_emergency_expense_prop_pt_est}%
    Berapakah estimasi titik untuk parameter tersebut?
\item\label{us_emergency_expense_prop_se_name}%
    Apa nama statistik yang kita gunakan untuk mengukur
    ketidakpastian estimasi titik?
\item\label{us_emergency_expense_prop_se_calc_w_pt_est}%
    Hitung nilai dari
    bagian~(\ref{us_emergency_expense_prop_se_name})
    untuk konteks ini.
\item
    Seorang komentator berita kabel berpendapat bahwa nilainya sebenarnya 50\%.
    Apakah ia seharusnya terkejut oleh data tersebut?
\item
    Misalkan nilai populasi yang sebenarnya diketahui sebesar 40\%.
    Jika kita menggunakan proporsi ini untuk menghitung ulang nilai pada
    bagian~(\ref{us_emergency_expense_prop_se_calc_w_pt_est})
    dengan menggunakan $p = 0.4$ sebagai pengganti $\hat{p}$,
    apakah nilai yang dihasilkan banyak berubah?
\end{parts}
}{}

\D{\newpage}

% 5

\eoce{\qt{Sampel air berulang\label{repeated_water_samples_prop}}
Sebuah organisasi nirlaba ingin memahami proporsi rumah tangga yang
memiliki kadar timbal tinggi dalam air minumnya.
Mereka memperkirakan sedikitnya 5\% rumah memiliki kadar timbal tinggi,
tetapi tidak lebih dari sekitar 30\%.
Mereka mengambil sampel acak 800 rumah dan bekerja sama dengan para pemilik
untuk memperoleh sampel air, lalu menghitung proporsi rumah tersebut
yang memiliki kadar timbal tinggi.
Mereka mengulangi proses ini 1.000 kali dan membentuk sebuah distribusi
proporsi sampel.
\begin{parts}
\item
    Apa nama distribusi ini?
\item
    Apakah Anda memperkirakan bentuk distribusi ini
    simetris, menceng kanan, atau menceng kiri?
    Jelaskan penalaran Anda.
\item
    Jika proporsi-proporsi tersebut tersebar di sekitar 8\%,
    berapakah variabilitas distribusinya?
\item
    Apa nama formal nilai yang Anda hitung pada bagian~(c)?
\item
    Misalkan anggaran para peneliti dikurangi sehingga mereka hanya
    mampu mengumpulkan 250 pengamatan per sampel, tetapi tetap dapat
    mengumpulkan 1.000 sampel.
    Mereka membentuk distribusi baru proporsi sampel.
    Bagaimana perbandingan variabilitas distribusi baru ini
    dengan variabilitas distribusi ketika setiap sampel
    memuat 800 pengamatan?
\end{parts}
}{}

% 6

\eoce{\qt{Sampel mahasiswa berulang\label{repeated_student_samples_prop}}
Di antara seluruh mahasiswa tahun pertama pada sebuah perguruan tinggi besar, 16\%
masuk daftar prestasi dekan pada tahun berjalan.
Sebagai bagian dari proyek kelas, para mahasiswa mengambil sampel acak 40 mahasiswa
dan memeriksa apakah mereka masuk daftar tersebut.
Mereka mengulangi proses ini 1.000 kali dan membentuk sebuah distribusi
proporsi sampel.
\begin{parts}
\item
    Apa nama distribusi ini?
\item
    Apakah Anda memperkirakan bentuk distribusi ini
    simetris, menceng kanan, atau menceng kiri?
    Jelaskan penalaran Anda.
\item
    Hitung variabilitas distribusi ini.
\item
    Apa nama formal nilai yang Anda hitung pada bagian~(c)?
\item
    Misalkan para mahasiswa memutuskan mengambil sampel lagi,
    kali ini sebanyak 90 mahasiswa per sampel,
    dan kembali mengumpulkan 1.000 sampel.
    Mereka membentuk distribusi baru proporsi sampel.
    Bagaimana perbandingan variabilitas distribusi baru ini
    dengan variabilitas distribusi ketika setiap sampel
    memuat 40 pengamatan?
\end{parts}
}{}
